% =============================================================================
% Chapter 7: Non-Functional Requirements
% =============================================================================

\chapter{Non-Functional Requirements}
\label{ch:nonfunctional}

\section{Reliability and Testing}

\begin{reqbox}{REQ-74}{\stable}
\textbf{Test-driven development.}\\
Every public method has its test written before its implementation. Pull
requests adding or modifying public API without a corresponding test are
rejected.
\end{reqbox}

\begin{reqbox}{REQ-75}{\stable}
\textbf{Coverage threshold.}\\
Test coverage measured with \code{pytest-cov} does not fall below 90\,\%
on the main branch. This is a hard CI gate, not an aspirational target.
\end{reqbox}

\begin{reqbox}{REQ-76}{\stable}
\textbf{Required edge-case tests.}\\
The test suite includes the cases enumerated in
Table~\ref{tab:required-edge-cases}, organized by family for audit ease.
\end{reqbox}

\begin{table}[H]
\centering
\caption{Required edge-case tests, grouped by area.}
\label{tab:required-edge-cases}
\small
\begin{tabularx}{\textwidth}{@{}lX@{}}
\toprule
\textbf{Area} & \textbf{Edge case}\\
\midrule
Load        & Empty VarFrame; single source file; dict-mode coordinate-tuple length mismatch; pandas DataFrame with non-string column names; NumPy array with mismatched \code{names}.\\
Pivot       & Datapoint mode with no subsequent pivot; \code{pivot()} on already-structured VarFrame; \code{auto\_detect=True} on ambiguous candidates; duplicate coordinates on every requested dimension raise \code{PivotDuplicateError}.\\
Diagnostics & 100\,\%-coverage VarFrame; all-NaN slice; single-point dimension.\\
Compute     & Invalid AST expression (syntax error, undefined variable); cyclic \code{.itceq} equations caught at parse time; non-differentiable NumPy with uncertainty guard (per-variable); \code{where} with \code{fill=None} on non-existent variable; \code{compute} with both symbolic and Monte Carlo methods on the same expression.\\
Ops         & \code{diff} with \code{window <= deg}; \code{interpolate} with non-monotonic axis-translation target; \code{interpolate} with \code{override=False} on existing coordinate; \code{expand} with duplicate existing dimension name; \code{itc.concat} with overlapping \code{along} values; \code{fitvalue} evaluation outside the original fit range producing \code{-1} tags; \code{combine} with mismatched grids; numerical operation on non-numeric dimension.\\
Compare     & Grid mismatch in \code{compare} without reference; multiple named inputs with overlapping deltas.\\
Integrate   & \code{np.nan} inside the integration domain with \code{skipna=False} (result NaN) and \code{skipna=True} (populated cells only, History records the fraction).\\
Immutability & In-place mutation of any stored array raises; \code{to\_numpy} returns a read-only view by default.\\
Reproducibility & Two identical load-and-operate runs yield identical state hashes; hashes are invariant to timestamps and user identity.\\
History     & Hash mismatch on reload; pipeline applied to incompatible VarFrame; comment preserved through \code{.itc} round-trip.\\
Modes       & Draft-mode export blocked without flag; mixed-mode binary operation rejected; \code{db.promote} / \code{db.demote} round trip preserves data.\\
Uncertainty & Mismatched uncertainty keys; correlation matrix non-symmetric; correlation matrix with $|r| > 1$; declared correlation between non-existent variables; repeat averaging reduces the random component by $1/\sqrt{N}$ and leaves the systematic component unchanged.\\
Axes        & \code{rotate} with unregistered axis; vector group cannot be resolved; non-orthogonal rotation matrix rejected; uncertainty propagated through rotation matches direct LPU.\\
\bottomrule
\end{tabularx}
\end{table}

\begin{reqbox}{REQ-77}{\stable}
\textbf{Property-based tests for math kernels.}\\
Property-based tests using Hypothesis~\citep{hypothesis} are required for:
\begin{itemize}[noitemsep,topsep=2pt]
  \item all AST operators (verifying analytical partial-derivative
    identities on randomized inputs);
  \item the propagation engine (verifying additivity of variances under
    independent inputs, correctness under known correlated inputs, and
    dimensional consistency);
  \item the \code{diff} kernel (verifying exact recovery of polynomials
    of degree $\le \code{deg}$);
  \item the rotation engine (verifying that
    $R^\top R = I$ is preserved, and that uncertainty propagated through
    rotation matches direct LPU computation).
\end{itemize}
\end{reqbox}

\section{Code Quality}

\begin{reqbox}{REQ-78}{\stable}
\textbf{Type hints.}\\
All public functions and methods carry full Python type hints.
\code{mypy --strict} is run in CI on the public API.
\end{reqbox}

\begin{reqbox}{REQ-79}{\stable}
\textbf{Docstrings.}\\
All public functions and classes carry NumPy-format docstrings with
Parameters, Returns, Raises, and at least one Examples section.
\end{reqbox}

\begin{reqbox}{REQ-80}{\stable}
\textbf{Linting and formatting.}\\
All code passes \code{ruff check} linting and \code{ruff format}
formatting. One tool covers both concerns. CI runs both checks on every
push and rejects non-conforming code.
\end{reqbox}

\begin{reqbox}{REQ-81}{\stable}
\textbf{Error messages.}\\
All exception messages include three pieces of information: the variable
or object involved; the operation attempted; and a suggested fix. Errors
derive from the \code{ITACAError} hierarchy
(Table~\ref{tab:exception-hierarchy}).
\end{reqbox}

\section{Architecture and Dependencies}

\begin{reqbox}{REQ-82}{\stable}
\textbf{NumPy-only core.}\\
Every \itc{} package depends only on NumPy and the standard library,
with two named exceptions: \code{io/} and \code{utils/}, which may
import pandas lazily as the DataFrame bridge (REQ-05, REQ-84). No
xarray, dask, or pandas imports appear in any other package, at module
level or inside a function.

The scope is stated as \emph{everything except two} rather than as a
list of the packages it covers. A list has to be extended whenever a
package is added, and a package added but not listed would inherit no
restriction at exactly the moment it is least reviewed; stated this
way, a new package is restricted by default and an exemption is a
deliberate act. At the 0.1.0 baseline the covered packages were
\code{core/}, \code{ops/}, and \code{uncertainty/}, which is how this
requirement was originally written; \code{pproc/} joined them at the M1
Phase B3b checkpoint without the text needing to change, which is the
property the restatement buys.

CI enforces this with a \code{ruff} import-policy rule, banning the
three packages repository-wide with per-file exemptions for the two
above. A guard test enforces it independently over the abstract syntax
tree, so the policy survives a lint-configuration regression. Both
sides discover packages rather than enumerate them, and the two
exemption declarations are checked against each other, so neither can
quietly stop covering a package the other still covers.
\end{reqbox}

\begin{reqbox}{REQ-83}{\stable}
\textbf{Pinned dependency version ranges.}\\
\itc{} declares an explicit PEP 440 specifier for every dependency in
\code{pyproject.toml}. Table~\ref{tab:dependency-versions} lists that
specifier and the role: where \code{pyproject.toml} declares the
dependency, or that it does not declare it yet.

The specifier is usually a bounded range, and it is not always one, so
this requirement says \emph{specifier} rather than \emph{range with an
upper bound}. Three of the DECLARED rows are deliberately otherwise.
\code{ruff} is an EXACT pin (see below). \code{setuptools} and
\code{setuptools-scm} carry a floor with no ceiling, because they are
build-time tools resolved by the frontend in an isolated environment,
where a ceiling pins a consumer's toolchain rather than this project's.
The two \code{external} rows carry no specifier at all and read
\code{system}, because they name system tools rather than Python
packages.

A package may be declared in more than one extra, and the role column
names one. \code{pandas} is the case: its role is \code{optional}, for
the extra of that name, and it is ALSO in \code{dev}, because a test
imports it at module scope, so the documented
\code{pip install -e ".[dev]"} fails at collection without it
(\code{ITACA-015}). That second membership is load-bearing and is
guarded by name in \code{tests/test\_packaging\_metadata.py}, since the
table cannot express it.

CI runs the full test suite against the LATEST compatible versions on
the lowest and highest supported interpreters, and against the declared
minimum SERIES of the runtime dependencies \code{numpy} and
\code{pandas} on the lowest. Three statements of what it does NOT do,
made here because this paragraph exists to replace an inferred promise
with a stated one, and an earlier draft of it overstated all three.
Interpreters between the ends of the supported range are not exercised:
\code{3.12} is supported and has no leg. The floor of every other row
is declared and not minimum-tested, which covers every development and
build-time dependency. And the minimum leg pins a series
(\code{1.26.*}, \code{2.1.*}) rather than the exact declared floor.
These are known limitations, not promises.

The table is CHECKED against \code{pyproject.toml} rather than
generated from it, in both directions, by
\code{tests/test\_packaging\_metadata.py}. Generated would lose the
\emph{role} and \emph{used for} columns, which are the reason this
belongs in a specification and not in a lockfile; checked gives the
same protection from drift. The check was written because the table had
drifted in both directions at once: it omitted six declared
dependencies, stated a range for \code{ruff} where the file carries an
exact pin, and carried five rows for packages declared in no extra at
all, so a reader could not tell a policy from a plan.

That last case is why the roles are split. A \code{planned} row is a
package this project intends to declare and has not; an
\code{external} row is not a Python package at all. Neither may name
something \code{pyproject.toml} declares, so a row cannot go on sitting
in an uncheckable role once the package becomes real.

The language baseline is normative in both directions: it may not fall
below a standard-library module the implementation already depends on.
The floor rose from 3.10 to 3.11 at the M1 Phase B3b checkpoint so that
\code{tomllib} is available, which is what lets the \code{.itceq}
parser (REQ-48) read the TOML structure of Section~\ref{sec:itceq} with
no dependency, no vendored reader, and no hand-written parser (OQ-28,
DD-32). The floor is declared once in \code{pyproject.toml} and
restated by the PyPI classifiers, the \code{ruff} target version, and
the lowest leg of the CI test matrix; the four are pinned to each
other, and to the standard-library modules the library imports, by
\code{tests/test\_python\_floor.py}.
\end{reqbox}

\begin{table}[H]
\centering
\caption{Dependency version policy for \itc{}. Bounds are PEP 440
specifiers for every row \code{pyproject.toml} declares; an
\emph{external} row names a system tool rather than a Python package
and carries \code{system}. Every declared row is checked against
\code{pyproject.toml} by
\code{tests/test\_packaging\_metadata.py}. The role says where that file
declares the dependency: \emph{required} in \code{dependencies},
\emph{optional} in an extra, \emph{dev} in the \code{dev} extra,
\emph{build} in \code{build-system.requires}, \emph{planned} for a
package this project intends to declare and has not, and
\emph{external} for what is not a Python package.}
\label{tab:dependency-versions}
\small
\begin{tabularx}{\textwidth}{@{}llXl@{}}
\toprule
\textbf{Dependency} & \textbf{Role} & \textbf{Used for} & \textbf{Range}\\
\midrule
\code{numpy}          & required  & Core arrays, AST, propagation, ops       & \code{>=1.26,<3.0}\\
\code{python}         & required  & Language baseline                        & \code{>=3.11,<3.14}\\
\code{pandas}         & optional  & DataFrame bridge                         & \code{>=2.1,<3.0}\\
\code{setuptools}     & build     & Build backend, PEP 639 metadata          & \code{>=77}\\
\code{setuptools-scm} & build     & Version derived from the repository       & \code{>=8}\\
\code{pytest}         & dev       & Test runner                              & \code{>=9.0.3,<10.0}\\
\code{pytest-cov}     & dev       & Coverage gate                            & \code{>=7.0,<8.0}\\
\code{hypothesis}     & dev       & Property-based testing                   & \code{>=6.100,<7.0}\\
\code{ruff}           & dev       & Linting and formatting                   & \code{==0.15.22}\\
\code{mypy}           & dev       & Static type checking                     & \code{>=1.10,<2.0}\\
\code{pre-commit}     & dev       & Local CI mirror                          & \code{>=3.7,<5.0}\\
\code{scipy}          & dev       & Geometry oracle only (DD-26)             & \code{>=1.11,<2.0}\\
\code{uncertainties}  & dev       & GUM oracle only (DD-25)                  & \code{>=3.2,<4.0}\\
\code{pyyaml}         & dev       & Release-gate checker reads workflow YAML & \code{>=6.0,<7.0}\\
\code{build}          & dev       & Release-integrity tests build artifacts  & \code{>=1.0,<2.0}\\
\code{matplotlib}     & planned   & Default plotting backend                 & \code{>=3.8,<4.0}\\
\code{plotly}         & planned   & Interactive plotting backend             & \code{>=5.20,<6.0}\\
\code{packaging}      & dev       & Version and specifier comparison in tests & \code{>=23.0,<26.0}\\
\code{smt}            & planned   & Surrogate modeling                       & \code{>=2.5,<3.0}\\
\code{tectonic}       & external  & LaTeX-backed PDF reports (preferred)     & system\\
\code{pdflatex}       & external  & LaTeX-backed PDF reports (fallback)      & system\\
\bottomrule
\end{tabularx}
\end{table}

The \code{ruff} pin is exact rather than a range, and that is
deliberate. Its authority is DD-29: \code{ruff} is not a library
\itc{} calls but a tool whose OUTPUT is a gate, and it is declared
twice, here and as a \code{pre-commit} hook revision that builds its
own isolated environment, so a range would let CI and the local hook
resolve different versions and run different rule sets. REQ-96 requires
the two to MIRROR each other and does not itself mandate an exact pin;
the pin is the means by which that mirror is kept, and
\code{tests/test\_tooling\_config.py} is what enforces it.

This paragraph is what DD-29 was waiting for. That decision recorded
itself as \emph{open, deferred to the author}, because REQ-83 then said
\emph{ranges} with an upper bound and an exact pin is neither, so by the
authority chain the SRS contradicted the pin. The wording above admits
an exact specifier for a tool of this kind, which is the amendment DD-29
named as one of its two exits. Whether that CLOSES the deferral is the
author's call and is not asserted here.

\code{scipy} and \code{uncertainties} are
development test oracles (DD-26, DD-25) and are barred from library
code by the import policy of REQ-82; their presence in this table is
not a relaxation of that rule.

\begin{reqbox}{REQ-84}{\stable}
\textbf{Optional dependency management.}\\
Optional dependencies (matplotlib, plotly, SMT, pandas, the LaTeX
toolchain) are imported lazily and checked at the point of use. A missing
optional dependency raises \code{MissingDependencyError} with an
actionable installation instruction (e.g.\
\code{pip install itaca[plot]}).

\emph{Optional} here is the USAGE sense, a dependency the library works
without, and it is deliberately wider than the \emph{optional} ROLE in
Table~\ref{tab:dependency-versions}, which means specifically
\emph{declared in an extra of \code{pyproject.toml}}. Of the five named
above, only \code{pandas} carries that role today; matplotlib, plotly
and SMT are \emph{planned} and the LaTeX toolchain is \emph{external},
because it is not a Python package. The \code{itaca[plot]} extra in the
example is likewise the shape the message will take, not an extra that
exists yet: it arrives with the plotting backend it names.
\end{reqbox}

\begin{reqbox}{REQ-85}{\stable}
\textbf{API parameter order.}\\
Canonical parameter order: \emph{(what, where, how)}: the variable or
object first, then location or dimension, then method or options.
Optional parameters beyond the first positional argument are
keyword-only.
\end{reqbox}

\begin{reqbox}{REQ-86}{\stable}
\textbf{English everywhere.}\\
All code, comments, docstrings, variable names, error messages,
\code{.itceq} files, and documentation are in English without exception.
\end{reqbox}

\begin{reqbox}{REQ-87}{\stable}
\textbf{American English with Z.}\\
\itc{} uses American English orthography consistently: \emph{organize}
(not organise), \emph{visualization} (not visualisation),
\emph{normalize} (not normalise). This reduces inconsistency and matches
the dominant convention in scientific Python documentation.
\end{reqbox}

\section{Performance and Memory}

\begin{reqbox}{REQ-88}{\stable}
\textbf{Vectorized operations.}\\
All operations on VarFrame arrays are vectorized via NumPy. Python loops
over array elements are not acceptable in performance-critical paths
(\code{compute}, propagation, \code{diff}, \code{smooth},
\code{interpolate}, \code{rotate}).
\end{reqbox}

\begin{reqbox}{REQ-89}{\stable}
\textbf{Memory transparency.}\\
\code{db.summary()} reports the in-memory footprint of a VarFrame.
Operations that potentially trigger large allocations (e.g.\ \code{expand}
broadcasting to a high-cardinality new dimension, dense pivot of a sparse
acquisition matrix) document the resulting size in their docstring and
print a warning when the projected allocation exceeds 1\,GB.
\end{reqbox}

\begin{reqbox}{REQ-90}{\stable}
\textbf{Sparse-aware handling of large datasets.}\\
For VarFrame instances whose coverage drops below a configurable
threshold (default 25\,\%), \itc{} provides sparse-aware variants of the
heavy operations (\code{compute}, \code{interpolate}, \code{integrate})
that operate only on populated cells. The sparse mode is opt-in via
\code{db.use\_sparse=True}; correctness must match the dense path
exactly, validated by paired tests. A common backing store (the underlying NumPy arrays plus a finite-mask) is used; \itc{} does not
introduce a separate sparse VarFrame type.
\end{reqbox}

\begin{reqbox}{REQ-91}{\stable}
\textbf{Lazy materialization of derived structures.}\\
The HistoryFrame, UncFrame, and CorrelationMatrix are materialized only
when first needed. A VarFrame loaded from raw data carries no UncFrame
until \code{set\_uncertainty} is called; no HistoryFrame until an
operation produces non-original values; no correlation matrix until
\code{set\_correlation} is called. This keeps the memory footprint of
exploratory workflows minimal.
\end{reqbox}

\section{Versioning and Release}

\begin{reqbox}{REQ-92}{\stable}
\textbf{Semantic versioning.}\\
\itc{} follows semantic versioning (\code{MAJOR.MINOR.PATCH}). Initial
public release is \code{0.1.0}. Breaking API changes increment
\code{MAJOR}; new backward-compatible features increment \code{MINOR};
bug fixes increment \code{PATCH}.

\textbf{Major version zero.} While \itc{} is at \code{0.y.z} it is in
initial development, and SemVer clause~4 applies: the public API is not
stable and a breaking change MAY ship in a \code{MINOR} bump. This is
stated rather than left implicit because \code{0.2.0} ships several of
them, and under the rule above with no carve-out that release would be
non-conforming; the requirement rather than the release is what was
wrong.

This requirement does NOT enumerate them, and states no count either.
The absence is a correction rather than an omission. It used to carry a
list. That list was already incomplete when the carve-out was written,
and the CHK-1 checkpoint then turned several further accepted inputs
into public refusals. A
second inventory inside a normative document goes stale independently
of the one it duplicates, and a bare count is the same artifact with
the evidence removed, so neither is kept here. The enumeration lives in
\code{CHANGELOG.md}, which the paragraph below already requires, and
that file is the single authority for what a release broke.

Every breaking change is marked BREAKING in \code{CHANGELOG.md} with
its migration, which is the obligation that replaces the version bump
while the major version is zero. The carve-out lapses at
\code{1.0.0}.
\end{reqbox}

\begin{reqbox}{REQ-93}{\stable}
\textbf{Conventional commits.}\\
All commit messages follow Conventional Commits. Allowed prefixes:
\code{feat:}, \code{fix:}, \code{docs:}, \code{test:}, \code{refactor:},
\code{chore:}, \code{perf:}, \code{build:}, \code{ci:}. CI rejects
non-conforming commits.
\end{reqbox}

\begin{reqbox}{REQ-94}{\stable}
\textbf{Changelog.}\\
\code{CHANGELOG.md} is updated with every version bump, following the
Keep a Changelog format~\citep{keepachangelog} (Added, Changed,
Deprecated, Removed, Fixed, Security).
\end{reqbox}

\begin{reqbox}{REQ-95}{\stable}
\textbf{Continuous integration.}\\
GitHub Actions runs on every push: \code{ruff check},
\code{ruff format --check}, \code{mypy --strict} (public API),
\code{pytest}, the coverage gate (threshold 90\,\%), and a build of this
specification. All checks must pass before a pull request may be merged.

The specification build compiles \code{docs/srs/main.tex} and fails on a
non-zero exit status OR on any error logged by the run, because a
nonstopmode build can emit a PDF while having logged errors, and a PDF
built from a document that fell into math mode is exactly the artifact
that looks like a success. It is listed here because a mandatory gate
absent from this enumeration is the defect this document's own 0.2.3
revision exists to close: two build defects reached \code{main} while
nothing compiled the document.
\end{reqbox}

\begin{reqbox}{REQ-96}{\stable}
\textbf{Pre-commit hooks aligned with CI.}\\
A \code{.pre-commit-config.yaml} mirrors the CI checks (\code{ruff check},
\code{ruff format --check}, \code{mypy}, fast \code{pytest -x}), running
locally on every commit. The local pre-commit configuration is a hard
requirement of the contribution workflow.

TWO DELIBERATE EXCEPTIONS to the mirror, and they are stated rather than
left as silent gaps between requirements.

FIRST, the specification build REQ-95 requires is NOT mirrored here. It
needs a TeX distribution, which is not a contribution prerequisite, and a
hook a contributor cannot run is a hook that gets skipped.

SECOND, since 2026-08-11, tests marked \code{guardproof} run in CI and NOT
in the local pre-push tier, which runs \code{-m "not guardproof"}. The
marker admits a test whose subject is a guard's own machinery, a mutation
companion proving a checker can still fail, rather than the library's
behavior. The author's decision answering BRF-076 is the authority: a push
must not be allowed to carry a change that breaks behavior, and proving
that a guard is well built answers a different question, which only
changes when the guard changes.

This exception is WEAKER than the first and the difference is the point.
The first moves a check nowhere local; the second moves it to a tier that
does not block the push.

WHAT COMPENSATES, STATED AT THE STRENGTH IT ACTUALLY HAS. A lane closing
its work runs \code{.claude/tools/closing\_ci\_check.py}, which refuses to
let the work be reported closed while CI is red, running or unknown
(INC-20260811-1745-itaca). That is an INSTRUCTION in the closing skill
with a mechanical ANSWER: the checker's verdict cannot be misread, because
it exits non-zero and prints the sentence to write instead. It is NOT an
enforcement that a session ran it, and nothing in this repository makes a
red CI block a push or a commit. An earlier draft of this paragraph said a
red CI "blocks the lane CLOSING"; a V and V review measured that claim as
stronger than the mechanism, and it is corrected rather than left to be
relied on.

The mirror claim in this requirement was itself proven incomplete by
ITC-20260802-2300, where three lanes ran a green local suite over a test
red on every CI run, so CI is the authority and the local tiers are a fast
pre-flight.
\end{reqbox}

\begin{reqbox}{REQ-97}{\draft}
\textbf{Citable releases.}\\
Tagged releases on GitHub are mirrored on Zenodo, generating a DOI for
each release. \itc{} is academically citable from \code{0.1.0}.
\end{reqbox}

\begin{reqbox}{REQ-108}{\draft}
\textbf{Published documentation site.}\\
\itc{} publishes a documentation site on GitHub Pages, built with
ProperDocs (the maintained MkDocs fork) and the material theme,
mirroring the pyflightstream site (DD-23). The navigation covers a
getting-started home page, the SRS chapters authored in markdown for
the site, \code{DECISIONS.md}, \code{OPEN\_QUESTIONS.md}, the
sister-library page, and an API reference. Following the pyflightstream
scheme, the reference and architecture pages are generated at build
time by \code{mkdocs-gen-files} and \code{mkdocs-literate-nav} from a
docs-build script that reads the live NumPy-format docstrings through
standard-library introspection, one page per public module and kept in
step with the milestone. The generator script and the mkdocs tooling
live in the docs build tree, are never imported by library code, and
nothing generated is committed, so the site cannot drift from the code
and the NumPy-only rule (REQ-82) is untouched. The renderer is a build
script rather than \code{mkdocstrings}, which pyflightstream dropped. A
\code{docs.yml} workflow builds the site with \code{--strict} on push
to \code{main} and deploys it to GitHub Pages. \code{properdocs},
\code{mkdocs-material}, \code{mkdocs-gen-files}, and
\code{mkdocs-literate-nav} are development-only dependencies, added to
the REQ-83 dependency table with version ranges and never runtime. The
published site does not replace the SRS as the authoritative
specification.
\end{reqbox}
